\section{Introduction}

  In previous work, we introduced a spectral entropy functional
  $S_\Pi[g] = \tfrac12 \log \det{}' A_g$,
  where $A_g$ is a Laplace-type elliptic operator associated with the
  projective substrate.
  In the infrared regime, the renormalized first variation of this functional
  was shown to generate an effective Einstein response, with higher-derivative,
  non-local, and anomaly contributions suppressed by powers of the
  pre-geometric scale $\ell_\chi$.
  This established the emergence of Einstein gravity as the dominant
  two-derivative sector of the spectral dynamics.

  However, the construction of Papers~1 and~2 was formulated in a
  Riemannian (elliptic) framework.
  While sufficient to derive the infrared field equations and their
  thermodynamic interpretation, this setting does not by itself address
  causal propagation.
  In particular, the elliptic operator $A_g$ encodes the spatial
  correlation structure of the substrate, but gravitational waves are
  dynamical, hyperbolic excitations propagating along null cones.

  The purpose of the present work is therefore to provide the
  Lorentzian completion of the spectral entropy dynamics.
  We define the Lorentzian action as
  \[
    S_\Pi^{(L)}[g]
    =
    \frac12 \Re \log \det{}' \mathcal{D}_g,
  \]
  equivalently via a Schwinger--Keldysh prescription,
  so that the quadratic kernel is real and retarded.
  The inverse operator entering the second variation
  $\delta^2 S_\Pi^{(L)}$ is thus the retarded Green operator $G_R$,
  ensuring causal propagation.

  We show that, in the regime
  $R\ell_\chi^2 \ll 1$,
  the quadratic operator derived from
  $\delta^2 S_\Pi^{(L)}$
  propagates exactly two transverse-traceless tensor modes of helicity
  $\pm 2$.
  Additional poles generated by higher-derivative corrections appear at
  $k^2 \sim \ell_\chi^{-2}$ and therefore lie outside the infrared
  window $\omega \ll \ell_\chi^{-1}$ relevant to present observations.
  In this domain, the theory reproduces the polarization content of
  general relativity.

  Beyond the leading Einstein sector, the spectral expansion generates a
  Weyl-squared contribution whose coefficient is fixed by the
  Seeley--DeWitt coefficient $a_4$.
  The Lorentzian signature of the background metric is not introduced as an
  independent postulate: it follows from the Born--Infeld admissibility
  constraints of the Cosmochrony framework, as established in~\cite{Beau2026q5b}.
  Linearization about Minkowski spacetime then yields a modified
  dispersion relation of the form
  \[
    \omega^2
    =
    c^2 k^2
    -
    \gamma\, \ell_\chi^2 k^4
    +
    \mathcal{O}(k^6),
  \]
  where $\gamma = 1/(180\zeta)$ and
  $\zeta = \mathcal{O}(1)$ encodes scheme-dependent normalization
  factors.
  In the natural spectral identification of the renormalization scale
  with the cutoff, $\mu=\ell_\chi^{-1}$, one has $\zeta=1/6$ and thus
  $\gamma=1/30$.

  This structure leads to a direct observational prediction.
  Gravitational-wave catalogues constrain modified dispersion relations
  of the form $\omega^2 = c^2 k^2 + A_4 k^4$.
  Comparison with the spectral result yields a bound
  \[
    \ell_\chi
    \lesssim
    \sqrt{180\,\zeta\,A_4^{\mathrm obs}},
  \]
  which becomes
  $\ell_\chi \lesssim \sqrt{30\,A_4^{\mathrm obs}}$
  in the natural scheme.
  The Lorentzian completion therefore provides the first direct
  gravitational-wave constraint on the pre-geometric scale introduced in
  the spectral framework.

  The structure of the paper is as follows.
  Section~\ref{sec:lorentzian_completion} defines the causal Lorentzian
  spectral action and its second variation.
  Section~\ref{sec:ir_spectrum} establishes the infrared propagating
  degrees of freedom and proves the transverse-traceless proposition.
  Section~\ref{sec:dispersion} derives the modified dispersion relation
  and its connection to the Seeley--DeWitt coefficients.
  Section~\ref{sec:gw_constraints} discusses the confrontation with
  gravitational-wave observations and the resulting bounds on
  $\ell_\chi$.
